\documentclass[10pt,a4paper]{amsart}
\usepackage[margin=19mm]{geometry}
\usepackage{amsmath,amssymb,mathtools}
\usepackage[scr=rsfs,cal=euler]{mathalfa}
\usepackage{xfrac}
\usepackage[unicode,hidelinks,bookmarks=false]{hyperref}
\newcommand{\ie}{i.e.\ }
\newcommand{\cf}{cf.\ }
\newcommand{\resp}{respectively\ }
\newcommand{\Subsection}{\subsection}
\providecommand{\nobreakdash}{\nobreak\hbox{-}}
\setlength{\emergencystretch}{2em}
\multlinegap0pt
\usepackage{bm}
\usepackage{bbm}
\usepackage{dsfont}
\usepackage{xspace}
\usepackage{cancel}
\usepackage{mathtools}
\usepackage{amsthm}
\makeatletter
\@ifundefined{thm}{\newtheorem{thm}{Theorem}}{}
\@ifundefined{lemma}{\newtheorem{lemma}[thm]{Lemma}}{}
\makeatother
\newcommand{\R}{\mathbb R}
\newcommand{\vect}{\text{Span}}
\title[Equiprojective polytopes in higher dimension]{Equiprojective polytopes in higher dimension}
\author{Alice Cousaert}
\date{Source: arXiv:2601.13095v1 (CC BY 4.0)}
\begin{document}
\maketitle

\noindent\emph{Source and licence.} Equiprojective polytopes in higher dimension. Version arXiv:2601.13095v1 (CC BY 4.0), distributed under the Creative Commons Attribution 4.0 International licence, as recorded in the arXiv licence metadata for this exact version and shown on the abstract page https://arxiv.org/abs/2601.13095v1. Full rights notice: licenses/arxiv-2601.13095-CC-BY-4.0.txt.\par\smallskip

\noindent\textit{Editorial scope.} Counted unit: the Lemma whose source label is \texttt{step\_3} in the article (source0.tex lines 722--725, source label \texttt{step\_3}), delivered together with the article's own complete original proof of it (source0.tex lines 727--765, one \texttt{proof} environment of 39 lines). No mathematics is written, repaired or re-derived: statement and proof are the article's own text line for line. Required context retained uncounted: balade (lines 604--607); step\_1 (lines 672--679); step\_2 (lines 693--696). Every definition, display and label of those lines is retained unchanged. Omitted from this package and not redistributed: 1--603, 608--671, 680--692, 697--721, 726--726, 766--1305. Nothing inside a retained excerpt is omitted or altered: every retained body is delivered exactly as the article's lines are written, so no omission receipt of any kind accompanies this package. Citations: the retained excerpts carry no citation command at all, so no bibliography entry of the article is transcribed and no reference is redistributed. Reference adaptation: none was needed and none was made; every reference inside the delivered unit resolves inside it, so no unresolved reference or citation is produced. Printed slips kept literal: no printed word, symbol, hypothesis or number of the counted statement or of its proof was altered. Layout and class: the article's own class is \texttt{article} with packages including babel, amsmath, amssymb, amsfonts, amsthm, array, calrsfs, enumitem, fancyhdr, geometry, graphicx, hyperref, inputenc, etoolbox; the wrapper is the lane's pinned \texttt{amsart} editorial wrapper at 10pt on A4, which restates the theorem-like environments the retained text needs and adds the article's own macro definitions that the retained text needs (\texttt{\textbackslash R}, \texttt{\textbackslash vect}), with extra wrapper packages (bm, bbm, dsfont, xspace, cancel, mathtools). The delivered numbering is the unit's own. Assets: no retained span contains a figure environment, a graphics-inclusion command, a PDF-page inclusion, a table or a TikZ picture; no image file is copied into this package, the package declares no asset dependency and no image file is redistributed with it. Licence: version arXiv:2601.13095v1 of this article is distributed under the Creative Commons Attribution 4.0 International licence, as recorded in the arXiv licence metadata for this exact version and shown on the abstract page https://arxiv.org/abs/2601.13095v1. A full rights notice is delivered at licenses/arxiv-2601.13095-CC-BY-4.0.txt. Characters: every retained excerpt is pure ASCII, so the delivered engine, XeLaTeX with its default Latin Modern text font, reports no missing character.
\begin{thm}[walking in the Grassmannian]
\label{balade}
    Let $P$ be a polytope of $\R^d$. We can go continuously from any admissible orthogonal to any other admissible orthogonal such that there are only finitely many degenerations along the path and at each of them, only one class of 2-faces is degenerating.
\end{thm}

\begin{lemma}[choosing the orthogonal of reference]
\label{step_1}
    By applying an isometry of $\R^d$ to the polytope $P$, we can get the following properties:
    \begin{itemize}
        \item No proscribed direction belongs to $e_1^\perp$. In this case, for each 2-faces $F$, $\vect[F]\cap e_1^\perp$ is a line spanned by some non-zero vector $\eta_F$.
        \item  The last coordinate of every $\eta_F$ is non-zero. We will assume that $\eta_{F,d}=1$
    \end{itemize}
\end{lemma}

\begin{lemma}
\label{step_2}
Let $P$ be a polytope satisfying the properties of Lemma \ref{step_1}, then we can go continuously from any admissible orthogonal to another included in $e_1^\perp$ while satisfying the properties of Theorem \ref{balade}.
\end{lemma}

\begin{lemma}
\label{step_3}
    Let $P$ be a polytope satisfying the properties of Lemma \ref{step_1}, then we can go continuously from any admissible orthogonal in $e_1^\perp$ to $W_r^\perp$ while satisfying the properties of Theorem \ref{balade}.
\end{lemma}

\begin{proof}
    Let $W_0^\perp\subset e_1^\perp$ be an admissible orthogonal which is not $W_r^\perp$. By Gaussian elimination, it admits a basis of the form \[ \begin{pmatrix}
        1& & & & & \\
         &\ddots& & &(0)& \\
        (0)& &1& & & \\
        x_2&...&x_{i-1}& & & \\
         & & &1& & \\
         &(0)& & &\ddots& \\
         & & & & &1
    \end{pmatrix} \] where each column represent a vector of the basis written in $(e_2,...,e_d)$. By the same argument as the perturbation $\gamma$ in the proof of Lemma \ref{step_2}, we can continuously go from $W_0^\perp$ to some perturbation $\tilde W_0^\perp$ while staying admissible along all the path, where $\tilde W_0^\perp$ has a basis \[ \begin{pmatrix}
        1& & & & & \\
         &\ddots& & &(0)& \\
        (0)& &1& & & \\
        x_2&...&x_{i-1}&\varepsilon& & \\
         & & &1& & \\
         &(0)& & &\ddots& \\
         & & & & &1
    \end{pmatrix} \] By Gaussian elimination, this operation allow the line of $x_j$ to go down one row. Iterating this process allows us to only consider basis looking like \[ \Gamma(x) := \begin{pmatrix}
        1& &(0) \\
         &\ddots& \\
        (0)& &1\\
        x_2&...&x_{d-1}\\
    \end{pmatrix} \]
    Then, we consider the path \[ t \in \mathopen [0,1  \mathclose] \mapsto \begin{pmatrix}
        1& &(0) \\
         &\ddots& \\
        (0)& &1\\
        (1-t)x_2&...&(1-t)x_{d-1}\\
    \end{pmatrix} \left(= \Gamma((1-t)x)\right) \] and we want to study the roots of the finitely many functions \[ D_F : t \in \mathopen [0,1  \mathclose] \mapsto \begin{vmatrix}
        1& &(0) &\eta_{F,2} \\
         &\ddots& & \vdots \\
        (0)& &1 & \eta_{F,d-1} \\
        (1-t)x_2&...&(1-t)x_{d-1} &1\\
    \end{vmatrix} \] After computation, we get that $\displaystyle D_F(t) = (1-t)\sum_{i=2}^{d-1}\eta_{F,i}x_i +1 $. They are affine functions and $D_F(1) = 1$ so they are non-zero. They admit at most one root in $\mathopen [0,1  \mathclose]$, and so there are only finitely many degenerations along this kind of path.
    
    If two non-parallel 2-faces are degenerating simultaneously, then \[ \varphi_F(x) = \sum_{i=2}^{d-1}\eta_{F,i}x_i = \sum_{i=2}^{d-1}\eta_{F',i}x_i = \varphi_{F'}(x) \] Where $\varphi_F$ and $\varphi_{F'}$ are linear forms in $x\in \R^{d-2}$. As before, we perturb $(x_i)_i$ in order to break this equality. We only need to find an $(\tilde x_i)_i$ such that $\varphi_F(\tilde x) \neq \varphi_{F'} (\tilde x)$, and then, with an $\varepsilon>0$ small enough, we are able to go continuously from $\vect(\Gamma(x))$ to $\vect(\Gamma(x+\varepsilon \tilde x))$ while being sure we do not add any new simultaneous degeneration. If there is no such $\tilde x$, then $\varphi_F = \varphi_{F'}$. Evaluating this relation in the canonical basis of $\R^{d-2}$ yields that $\eta_{F,i} = \eta_{F',i}$ for all coordinate. Thus $\eta_F = \eta_{F'}$, which is a contradiction. So there exists an $\tilde x$ such that $\varphi_F(\tilde x) \neq \varphi_{F'} (\tilde x)$.
    
    Then, by concatenating the finitely many needed perturbations and the path $t\mapsto \Gamma((1-t)x)$ for some $x$, we get the path we wanted.
\end{proof}

\end{document}
